\chapter{研究背景}
\chaptereng{Research Background}

\section{研究背景}

燃烧器类型与结构设计演进，平面燃烧器早期用于火焰高度、速度等基础参数测量及光学技术校准，凭借简单几何结构便于直接/间接测量火焰特性。之后采用多孔陶瓷结构（如15mm直径圆形多孔表面、0.2mm通道），搭配液体热稳定系统，实现火焰稳定控制；通过坐标变换等数值方法优化，适配详细反应机理模拟。多孔板或者多喷嘴阵列燃烧器通过多喷嘴或孔道设计强化燃料与氧化剂混合，降低回火风险，适配氢富集燃料。其中狭缝式多孔板：狭缝宽度0.32-1.0mm，固定长度6mm，通过调整孔径控制火焰形态，发现狭缝宽度减小会促进壁面点火替代流体动力回火。微尺度多喷嘴阵列：24个喷嘴，四喷嘴为一单元形成内回流区，壁面附近形成外回流区，通过高速射流诱导烟气再循环实现稳定燃烧，喷嘴出口速度可达40-100m/s。喷射稳定燃烧器采用同轴射流或热同流设计，如乙醇压力旋流喷雾燃烧器，通过二次燃烧器生成热稀释同流，适配温和燃烧工况，优化雾化机制，在热稀释环境中实现液体射流快速破碎，减少高温峰值（$<$2000K），降低NO$_x$排放\cite{xu2018characterization}。

\begin{figure}[H]
	\centering
	\begin{subfigure}{0.45\textwidth}
		\centering
		\includegraphics[width=\linewidth]{image4.png}
		\label{fig:initial_burner}
	\end{subfigure}
	\hfill
	\begin{subfigure}{0.45\textwidth}
		\centering
		\includegraphics[width=\linewidth]{image5.png}
		\label{fig:center_burner}
	\end{subfigure}
	\caption{初始平面火焰燃烧器与中心燃烧器}
	\label{fig:burners}
\end{figure}

燃烧稳定性与机制优化方面，对于回流区调控：通过内/外回流区耦合多喷嘴阵列、射流诱导烟气再循环，形成高温点火核心，拓宽稳定燃烧范围。对于不稳定性抑制：针对平面燃烧器的扩散热振荡，通过数值模拟优化反应机理，明确临界振荡条件与脉动周期；多喷嘴设计通过火焰相互作用抑制自激振荡\cite{liu2022investigation}。对于回火防护：氢富集燃料场景中，采用无旋高速喷嘴、多孔板狭缝结构，避免火焰逆流传播，同时通过壁面温度控制降低自动点火风险。燃烧模式的创新扩散热振荡控制：首次在常压下观测到甲烷空气火焰的扩散热脉动，通过高速成像、光谱分析明确振荡周期与自由基脉动的关联性。温和燃烧适配：通过热稀释同流实现无焰燃烧，喷雾燃烧中弱化内火焰面，降低温度梯度与污染物生成\cite{pers2024transition}。

低排放与燃料适应性发展方面，低NO$_x$技术通过缩短烟气停留时间、降低火焰峰值温度，实现NO$_x$排放$<$20ppm；利用热损失调控进一步降低高温区贡献。在氢富集燃料适配方面，针对氢高火焰速度、低淬熄距离特性，优化多喷嘴间距、狭缝结构，抑制回火\cite{koser2017investigation}；微尺度阵列燃烧器通过高速射流设计，适配氢含量0-100\%的混合燃料。同时有些平面火焰燃烧器拓展至煤粉颗粒燃烧，通过OH-PLIF诊断捕捉挥发分燃烧与焦炭氧化阶段，适配63-200$\mu$m粒径范围\cite{koser2019multi}。

\section{国内外研究现状}

\subsection{燃烧器的发展}

\subsubsection{平面火焰燃烧器、中心燃烧器的发展}

燃烧器类型与结构设计演进，平面燃烧器早期用于火焰高度、速度等基础参数测量及光学技术校准，凭借简单几何结构便于直接/间接测量火焰特性。之后采用多孔陶瓷结构（如15mm直径圆形多孔表面、0.2mm通道），搭配液体热稳定系统，实现火焰稳定控制；通过坐标变换等数值方法优化，适配详细反应机理模拟。多孔板或者多喷嘴阵列燃烧器通过多喷嘴或孔道设计强化燃料与氧化剂混合，降低回火风险，适配氢富集燃料。其中狭缝式多孔板：狭缝宽度0.32-1.0mm，固定长度6mm，通过调整孔径控制火焰形态，发现狭缝宽度减小会促进壁面点火替代流体动力回火。微尺度多喷嘴阵列：24个喷嘴，四喷嘴为一单元形成内回流区，壁面附近形成外回流区，通过高速射流诱导烟气再循环实现稳定燃烧，喷嘴出口速度可达40-100m/s。喷射稳定燃烧器采用同轴射流或热同流设计，如乙醇压力旋流喷雾燃烧器，通过二次燃烧器生成热稀释同流，适配温和燃烧工况，优化雾化机制，在热稀释环境中实现液体射流快速破碎，减少高温峰值（$<$2000K），降低NO$_x$排放\cite{xu2018characterization}。

燃烧稳定性与机制优化方面，对于回流区调控：通过内/外回流区耦合多喷嘴阵列、射流诱导烟气再循环，形成高温点火核心，拓宽稳定燃烧范围。对于不稳定性抑制：针对平面燃烧器的扩散热振荡，通过数值模拟优化反应机理，明确临界振荡条件与脉动周期；多喷嘴设计通过火焰相互作用抑制自激振荡\cite{liu2022investigation}。对于回火防护：氢富集燃料场景中，采用无旋高速喷嘴、多孔板狭缝结构，避免火焰逆流传播，同时通过壁面温度控制降低自动点火风险。燃烧模式的创新扩散热振荡控制：首次在常压下观测到甲烷空气火焰的扩散热脉动，通过高速成像、光谱分析明确振荡周期与自由基脉动的关联性。温和燃烧适配：通过热稀释同流实现无焰燃烧，喷雾燃烧中弱化内火焰面，降低温度梯度与污染物生成\cite{pers2024transition}。

低排放与燃料适应性发展方面，低NO$_x$技术通过缩短烟气停留时间、降低火焰峰值温度，实现NO$_x$排放$<$20ppm；利用热损失调控进一步降低高温区贡献。在氢富集燃料适配方面，针对氢高火焰速度、低淬熄距离特性，优化多喷嘴间距、狭缝结构，抑制回火\cite{koser2017investigation}；微尺度阵列燃烧器通过高速射流设计，适配氢含量0-100\%的混合燃料。同时有些平面火焰燃烧器拓展至煤粉颗粒燃烧，通过OH-PLIF诊断捕捉挥发分燃烧与焦炭氧化阶段，适配63-200$\mu$m粒径范围\cite{koser2019multi}。

\subsubsection{燃气轮机的发展}

在结构设计优化方面，非预混平面绝热火焰炉围绕强化掺混均匀性、提升工况适应性、优化测量便利性持续演进。早期简单喷嘴布局适用于低氢含量燃料基础测试，NASA的GRC（N1）型微混燃烧器通过交叉射流实现快速掺混，为多喷嘴阵列设计奠定基础\cite{cabra2004turbulent}。针对富氢燃料回火问题，GE公司的多管混合（MTmixer）喷嘴密集布置毫米级微混合管，利用高流速抑制回火，单筒全尺寸实验NO$_x$排放低于10ppm\cite{li2023motion}。国内方面，华北电力大学设计3$\times$3阵列微混燃烧器，采用燃料轴向进气、空气径向双排布孔结构，在60\%H$_2$+40\%CH$_4$工况下稳定燃烧\cite{chen2023富氢}；中国科学院工程热物理所提出中心1个加外圈6个的7喷嘴阵列，氢含量60\%时NO$_x$排放不高于10ppm，丰富了设计体系\cite{liu2021含氢}。早期非预混平面绝热火焰炉仅适用于常温常压、低负荷工况，随着燃气轮机向F级、H级升级，GE、三菱日立等厂商通过改进加热与压力控制模块，实现600-800K、1-2MPa严苛工况模拟。GE在先进燃氢透平项目中，利用非预混平面绝热火焰炉完成氢含量90\%以上燃料超100小时实验。国内研究适配减排需求，华北电力大学装置通过精准控流，实现喷嘴出口流速20-30m/s、热功率22.5-77.7kW宽范围调节；中国科学院工程热物理所优化壁面冷却结构，达成1920K绝热火焰温度下长时间稳定实验，为污染物排放研究提供支撑\cite{correia2015spray, venkataraman2011class}.

\begin{figure}[H]
	\centering
	\includegraphics[width=0.7\textwidth]{image6.png}
	\caption{微预混燃气轮机火焰流场\cite{cabra2004turbulent}}
	\label{fig:micromix_flow}
\end{figure}

测量技术集成创新贯穿非预混平面绝热火焰炉发展的全程，从常规参数测量迈向精细化表征。火焰与污染物测量中，PLIF技术与烟气分析仪结合实现同步测量，OH-PLIF技术表征火焰反应区，动态压力传感器捕捉压力脉动，为热声振荡研究提供支撑，大幅提升实验能力\cite{york2013development}。

同时在富氢燃烧研究中，非预混平面绝热火焰炉作用关键。掺混均匀性研究中，华北电力大学证实空气径向射流加燃料轴向射流掺混方式均匀性指数达0.978，中国科学院工程热物理所明确预混段长度35mm为临界值；稳焰机制研究中，GE揭示局部逆压梯度稳焰机制，三菱日立验证多孔喷射结构抑回火效果，国内明确富氢微混火焰依赖剪切层拉伸效应稳焰\cite{asai2014multiple, zhang2015temperature}；低NO$_x$排放优化中，GE、中国科学院工程热物理所分别实现氢含量60\%时NO$_x$排放低于25ppm、不高于10ppm，为减排提供数据支撑\cite{liu2021cold}.

\subsubsection{非预混火焰炉的测量发展}

由于之前发展的燃烧器不能达到足够高的火焰温度，容易不稳定和烧蚀，因此要设计新型的非预混平面绝热火焰炉。非预混火焰炉的测量方法核心围绕流场、温度、物种浓度、燃烧稳定性等四个关键参数，以非侵入式光学诊断为主。

对于流场特性的测量，瑞利散射成为火焰温度及组分浓度测量的重要手段，为燃烧过程诊断提供数据支撑。用Hencken炉平面火焰对滤波瑞利散射信号比-温度曲线进行精准标定，成功实现10 kHz高频二维测温，清晰揭示了单/双旋流火焰从V形向M形、从合并态向分层态的形态转变规律，该研究不仅验证了滤波瑞利散射技术在旋流燃烧诊断中的适配性，也为Hencken炉标定技术的应用提供了参考\cite{tao2020measurements}。在替代燃料燃烧诊断领域，将甲烷火焰的拉曼/瑞利散射测量系统拓展至二甲醚火焰，通过与CO平面激光诱导荧光技术联用，实现了H$_2$、O$_2$、N$_2$、CO等多组分浓度与温度的同步获取，为替代燃料燃烧过程的多参数综合诊断提供了可行方案\cite{li2024investigation, barlow2009raman}。

温度测量方面，相干反斯托克斯拉曼光谱（CARS）在高精度测量火焰温度方面具有优势，尤其适用于高压力、高温度或有干扰的场景\cite{shi2022compact}。在绝热火焰温度校准、爆震通道内高温测量（1300-2300K）、煤颗粒燃烧环境温度表征等方面应用广泛。基于氮气（N$_2$）振动拉曼散射，单脉冲精度可达2\%-4\%；采用飞秒/皮秒混合脉冲技术，可抑制非共振背景干扰，适配高压（1-10bar）环境。双颜色OH平面激光诱导荧光在二维温度场测量中兼顾空间分辨率与温度敏感性\cite{shi2022compact, barlow2009raman}。通过探测OH自由基的两个温度敏感跃迁，利用荧光强度比反演温度；空间分辨率达40-50$\mu$m，精度4\%-7.7\%\cite{hsu2021khz}。热电偶测量用于点式温度校准，配合辐射修正获取真实气体温度。

物种浓度与燃烧产物测量方面，平面激光诱导荧光可用于可视化自由基（OH、CH、HCO）或中间产物（CH$_2$O）的空间分布，识别反应区位置。在非预混火焰的反应区定位、煤颗粒燃烧的挥发分/焦炭燃烧阶段区分、爆震火焰的OH分布表征有广泛应用。气相色谱用于定量分析燃烧产物中的主要物种（CO、CO$_2$、H$_2$、未燃燃料）浓度，通过石英探针取样（针尖直径100$\mu$m），结合色谱柱分离，测量不确定度$\pm$0.2\%-0.9\%\cite{wagner2012absolute}。

燃烧稳定性与动态特性测量方面，高速摄影与化学发光成像采用高速CMOS相机，记录CH、C$_2$等化学发光信号，可识别火焰振荡周期（如7ms）、旋转螺旋结构等不稳定现象。激光散射与透射测量采用激光多普勒测速（LDV）或相位多普勒粒子分析（PDA），获取粒子速度与粒径信息，适用于稀相燃烧系统\cite{liu2018two}.

\begin{table}[H]
	\centering
	\caption{燃烧器测量方法和主要成果发展汇总}
	\label{tab:measurement_methods}
	\begin{tabular}{p{5cm}p{10cm}}
		\toprule
		测量方法 & 主要成果 \\
		\midrule
		氪X射线荧光光谱 & 量化多元素扩散燃烧器近表面（0.1-0.5 mm）燃料/空气快速混合过程；获得火焰温度分布\cite{alqasemi2023temperature} \\
		混合飞秒/皮秒振动相干反斯托克斯拉曼光谱 & 实现1300-2300 K、1-10 bar高压层流火焰温度测量\cite{braun2022spatially} \\
		双颜色OH平面激光诱导荧光 & 实现爆震通道内二维温度场测量，精度4\%-7.7\%\cite{hsu2021khz} \\
		气相色谱（GC） & 实现臭氧激活二甲醚低温冷火焰（779-885 K）稳定与表征；测定0.4-1.4当量比下冷火焰传播速度 \\
		高速OH平面激光诱导荧光（OH-PLIF）、相干反斯托克斯拉曼光谱（CARS） & 10 kHz帧率捕捉单煤颗粒燃烧OH分布，区分挥发分与焦炭燃烧阶段；量化富氧环境下煤颗粒燃烧边界层特性\cite{koser2019multi, shi2022compact} \\
		\bottomrule
	\end{tabular}
\end{table}

\subsection{TDLAS测量与火焰标定发展现状}

\subsubsection{TDLAS技术发展现状}

可调谐二极管激光吸收光谱（TDLAS）技术作为一种非侵入式光学测量手段，凭借高灵敏度、快速响应、物种特异性强等优势，已成为燃烧诊断、工业过程监测、环境检测等领域的核心技术支撑。随着现代工业向高精度、高效率、低排放方向转型以及极端工况下测量需求的不断增加，TDLAS技术逐步朝着单谱线测量、多参数协同、宽波段适配、智能化反演的方向演进，有效解决了传统测量方法中侵入式干扰、响应滞后、多组分耦合解算困难等关键问题\cite{liu2022tunable}。

TDLAS技术应用于火焰标定测量的基本原理是比尔-朗伯（Beer-Lambert）定律，当频率为$\nu$的单色激光穿过气体介质后，若光子能量与待测气体组分的吸收峰对应频率重合，部分激光将会被气体吸收，激光光强的透过率$T_\nu$可用沿光程的吸收率积分表示：
\begin{equation}
	T_\nu = \left(\frac{I}{I_0}\right)_\nu = \exp\left[-\int_0^L \alpha(l,\nu) \cdot dl\right]
	\label{eq:beer_lambert}
\end{equation}
式中：$I_0$和$I$分别为入射激光光强和透射激光光强；$\alpha(l,\nu)$为待测体系在吸收路径上$l$处频率为$\nu$的激光的吸光度；$L$为待测体系吸收距离。

定义吸光度$\alpha(l,\nu)$为：
\begin{equation}
	\alpha(l,\nu) = \sum_j \chi(l) \cdot S_j[T(l)] \cdot f_j[(\nu - \nu_{0,j}), T(l)]
	\label{eq:absorbance}
\end{equation}
式中：$f_j$是跃迁$j$的线型函数；$S_j$是跃迁$j$的谱线强度；$\nu_{0,j}$是跃迁$j$的中心跃迁频率；$\chi(l)$是待测体系$l$处的物种浓度。

温度$T$和线强$S_j(T)$存在函数关系：
\begin{equation}
	S_j(T) = S_j(T_0) \frac{Q_j(T_0)}{Q_j(T)} \frac{T_0}{T} \exp\left[-\frac{hc E_j''}{k}\left(\frac{1}{T} - \frac{1}{T_0}\right)\right] \times \left\{\frac{1-\exp(-hc\nu_{0,j}/kT)}{1-\exp(-hc\nu_{0,j}/kT_0)}\right\}
	\label{eq:line_strength}
\end{equation}
式中：$h$为普朗克常数；$k$为玻尔兹曼常数；$c$为光速；$T_0$是标准状况温度；$E_j''$为吸收跃迁的低能级能量；$Q(T)$是配分函数。

早期TDLAS测量技术多依赖多谱线组合方案，例如通过选取5条低态能级分布均匀的H$_2$O近红外吸收线，采用玻尔兹曼图法可实现1000-1600K温度范围的测量，相对误差控制在1\%以内，且在超燃冲压发动机扩张段的实测中展现出良好稳定性，显著提升了对光谱参数误差和环境干扰的容忍度。随着单谱线测量技术的突破，其核心优势在于简化光学系统设计、降低多谱线耦合干扰。基于模糊PID控制器的单谱线迭代测量方法，利用Voigt线型的微分方程可从单一吸收谱线中同时解耦温度、浓度和压力参数，在300-2300K范围内实现了2.5\%以内的相对温度误差，在水温恒温箱实验中温度误差控制在5$^\circ$C以下，并成功应用于对冲火焰的二维温度与水汽浓度成像\cite{hsu2021khz, liu2022tunable}。该方法通过迭代算法有效分离重叠谱线，避免了传统多谱线方案中基线提取复杂、系统集成繁琐的问题。此外，基于单一量子级联激光器的多组分测量方案也得到发展，利用单激光器覆盖CO、N$_2$O、H$_2$O的特征吸收波段，实现了大气多组分的同时测量，进一步拓展了单谱线技术的应用场景。

多参数协同测量与宽波段适配技术的发展显著提升了TDLAS技术的工况适应性。在燃烧诊断领域，光纤耦合的近红外/中红外（NIR/MIR）多波长TDLAS系统被成功开发，该系统集成4支激光器覆盖1391.7 nm、1469.3 nm（H$_2$O）、4854.6 nm（CO）和4556.2 nm（CO$_2$）波段，可实现温度、CO、CO$_2$、H$_2$O的同时测量，温度测量误差小于5\%，浓度测量误差优于10\%\cite{shi2022compact}。在复杂燃烧环境中，改进型直接吸收光谱（DAS）方法通过波长分复用技术分离OH自由基与H$_2$O的吸收信号，利用1502.69 nm（OH）和1349.23-1349.27 nm（H$_2$O）的谱线组合，实现了甲烷/空气预混火焰中OH自由基浓度、H$_2$O浓度及温度的同步测量，测量结果与模拟值高度吻合\cite{wen2022concentration}。

智能化反演算法与tomography技术的结合推动TDLAS技术向二维/三维场分布测量方向拓展。基于堆叠长短期记忆网络（LSTM）的质量分层温度成像网络（QHT-LSTM），通过32束激光阵列的视线测量数据，可同时输出粗精度（实时监测，temporal resolution $>$27 kfps）和高精度（离线分析）温度图像，在SNR=35 dB条件下，较传统CNN模型的重建误差降低43.62\%-57.80\%\cite{si2022quality}。TDLAS技术的工程化应用不断突破极端工况限制，在超燃冲压发动机领域，多波长TDLAS系统通过16路光路布置实现了扩张段横截面的温度测量，有效覆盖了点火、燃烧、熄火全流程的宽温度范围\cite{agarwal2023investigations, wagner2012insitu}。在工业燃烧设备中，基于CO$_2$在4.2$\mu$m波段的高温敏感吸收峰，采用TDLAS技术可测量Hencken平面火焰的温度与CO$_2$体积分数二维分布，在空气流量55 L/min、当量比0.7-0.9工况下，火焰温度接近绝热工况，为激光诊断技术提供了可靠的标定基准\cite{liu2018two}。

TDLAS用于火焰标定是依托火焰的可控实验条件（如固定燃料类型、压力、当量比等），建立激光吸收测量值与火焰真实参数的定量关联。通过测量值与真实值的对比与修正，校准TDLAS系统的光谱解读精度，修正参数偏差、提高光谱模型适用性，最终实现对火焰的精准定量测量。其中火焰温度是核心标定参数之一，包括火焰局部温度和温度分布等，同时也要标定不同燃烧产物在火焰中的绝对浓度及空间分布，这也间接标定TDLAS系统的光谱拟合算法、吸收线参数（如低态能量、碰撞展宽系数）的准确性，验证光谱模型在火焰强温度梯度、多组分干扰下的适用性\cite{zhou2005development}。

在低压均匀火焰的温度标定方面，设计了基于扫描波长调制光谱的TDLAS传感器，控制火焰压力、燃料类型等关键工况，将TDLAS测量结果与理论燃烧模型计算值对比，建立测量值与真实值的关联框架，为TDLAS在火焰环境中的准确性验证提供了实验方案，明确了光谱拟合算法在低压均匀火焰中的有效性。在非预混层流对冲扩散火焰的物种浓度标定上，实现了乙炔（C$_2$H$_2$）浓度的TDLAS标定与原位测量，通过多位置扫描获取火焰中C$_2$H$_2$的浓度分布，结合燃烧化学动力学模型，验证TDLAS在强温度梯度环境中的测量稳定性，为多组分、强梯度火焰的TDLAS应用提供了标定参考\cite{wagner2012absolute, wagner2012insitu}。在高温火焰的温度精准标定方面，基于2.9$\mu$m近红外TDLAS传感器，设计H$_2$O双色测温方案，实现了高温火焰的温度标定。选择H$_2$O具有不同低态能量的特征吸收线，利用其温度敏感性差异构建双色测温模型，抵御其他燃烧产物的吸收干扰，建立了高温火焰温度的多源验证标定方法，明确了TDLAS在高温环境下的测温精度与适用边界\cite{li2011h2o, zhang2018reconstruction, wang2023吸收}.

\begin{figure}[H]
	\centering
	\includegraphics[width=0.7\textwidth]{image7.png}
	\caption{TDLAS系统用于超燃冲压发动机测温示意图\cite{huang2023基于}}
	\label{fig:tdlas_scramjet}
\end{figure}

\begin{figure}[H]
	\centering
	\includegraphics[width=0.7\textwidth]{image8.png}
	\caption{TDLAS测量平面火焰炉实验装置\cite{wang2023吸收}}
	\label{fig:tdlas_setup}
\end{figure}

\subsubsection{火焰标定技术发展现状}

火焰标定技术作为燃烧温度测量领域的核心基准支撑技术，其核心价值在于为燃烧过程温度监测提供可溯源、高精度的量值依据，是保障工业制造、能源利用、环境保护等领域温度测量准确性的关键。随着现代工业向高效、低排放、高精度方向转型，燃烧温度测量对标定技术的要求日益严苛，不仅需要满足常规工况下的精度需求，更需适配工业炉窑、燃气轮机、核退役等复杂恶劣环境的现场标定场景\cite{smith2014scanned, wagner2012insitu}。国际温标（ITS-90）作为火焰标定技术的核心溯源基准，其精准实现与传递是该领域的核心研究目标，推动火焰标定技术逐步从实验室固定装置向便携化、原位化、智能化方向演进，以解决传统技术在现场应用中量值溯源困难、测量不确定性高、工况适配性差等突出问题\cite{huang2023基于}.

早期火焰标定技术多依赖大型实验室固定装置，存在体积庞大、操作复杂、无法现场部署等局限，难以满足工业现场实时校准需求。为突破这一瓶颈，欧洲EURAMET组织通过EMPIR计划支持的EMPRESS系列项目（14IND04、17IND04）开展了系统性研究，推动火焰标定技术发展。该系列项目由英国国家物理实验室（NPL）牵头，联合丹麦技术大学（DTU）、马德里卡洛斯三世大学（UC3M）等多家科研机构、高校及企业，共同开展研发，其成果推动了当前火焰标定技术的发展。其中，由英国国家物理实验室开发的便携式标准火焰系统作为这个项目的标志性成果，改变了传统火焰标定的应用场景限制，为现场溯源提供了一定的可行方案\cite{morley2024empir, jones2021empir, strathclyde2021empress2}.

标准火焰装置的设计与性能优化是火焰标定技术发展的核心载体。英国国家物理实验室开发的便携式标准火焰基于微扩散燃烧器，采用丙烷/空气作为燃料，通过精准控制当量比（$\phi$=0.8-1.4），实现了2040-2260K的稳定温度输出。该装置的核心技术创新体现在多方面：采用流量不确定度小于0.5\%的高精度质量流量计精准控制燃料与空气配比，结合氮气协流技术有效抑制火焰扰动，使燃烧器中心$\pm$10mm区域内温度均匀性显著提升；通过Rayleigh散射测温技术实现对ITS-90的直接溯源，综合测量不确定度仅为0.5\%（k=1），长期重复性优于0.2\%，性能远超传统标定装置\cite{sutton2016npl}。采用模块化设计实现便携化，可在3小时内完成现场部署与校准，已在欧洲多国工业现场与实验室循环使用，成功应用于FTIR光谱仪、高光谱成像系统等燃烧测温设备的准确性验证。其温度场特性经实验验证，在燃烧高度20mm处，不同当量比下的温度波动小于$\pm$3K，为红外光谱分析测量技术提供了理想的标定环境\cite{npl2025optical, jones2018empress}.

标准火焰的温度场特性与不确定性控制是火焰标定技术的核心研究内容。英国国家物理实验室通过为期10天的连续测试验证了该标准火焰的长期稳定性，结果表明，在20mm燃烧高度处，不同当量比下的火焰温度变化率均小于0.2\%，且与进气温度、大气压力等环境因素无明显相关性\cite{morley2024empir}。温度场分布测试显示，有无氮气协流条件下，火焰在水平方向$\pm$10mm范围内温度均匀性良好。不确定性预算分析表明，摩尔折射率数据、流量计精度、化学平衡假设、激光稳定性等是影响标定精度的主要因素，通过优化设计与误差补偿，最终实现了0.45\%-0.5\%的综合不确定度\cite{euramet2018enhancing}。此外，标准火焰的物种成分已通过GASEQ化学平衡求解器精准表征，明确了不同当量比下CO$_2$、CO、H$_2$O等主要燃烧产物的摩尔分数分布和温度变化，为基于光谱法的温度测量仪器校准提供了基础数据支撑，也为后续在进行标定时的当量比调节提供了参考\cite{flanagan2023new}.

\begin{figure}[H]
	\centering
	\includegraphics[width=0.7\textwidth]{image9.png}
	\caption{有无氮气协流温度的均匀性\cite{fateev2019validation}}
	\label{fig:n2_coflow}
\end{figure}

\begin{figure}[H]
	\centering
	\includegraphics[width=0.7\textwidth]{image10.png}
	\caption{不同当量比温度的变化\cite{fateev2019validation}}
	\label{fig:temp_vs_phi}
\end{figure}

基于标准火焰的标定方法创新是火焰标定技术发展的重要方向。红外发射光谱标定技术凭借非侵入性优势得到广泛应用，其与标准火焰的结合形成了高效的标定体系\cite{jones2021work, euramet2018publishable}。DTU开发的单线光谱仪与UC3M研发的中红外（2-5.5$\mu$m）高光谱成像系统，通过测量火焰中CO$_2$、CO和H$_2$O的特征发射谱，结合HITEMP-2010光谱数据库进行模拟对比，实现了火焰温度的精准反演。实验验证表明，该方法测量结果与NPL标准火焰的Rayleigh散射测量结果一致性优于1\%，在$\phi$=1.0工况下，中心区域温度测量偏差小于0.23\%\cite{euramet2021final}。此外，机器学习技术的引入为火焰温度反演提供了新路径，通过1800-2500cm$^{-1}$和3000-4200cm$^{-1}$双波段光谱数据反演温度分布，预测误差小于0.5\%，在垃圾焚烧炉等工业场景中已得到初步应用\cite{zhou2003development}.

当前火焰标定技术虽已取得显著进展，但仍面临诸多挑战。在极端工况适配方面，超高温（$>$2260K）、强辐射、高速瞬态等特殊条件下的标定方法尚不完善，现有标准火焰的温度范围和环境适应性仍需拓展。

\section{本文研究内容}

综合以上文献分析，现有研究存在以下不足：在燃烧器结构方面，现有平面炉缺乏针对3000K超高温工况的主动热管理设计，传统插管结构存在热膨胀失配和烧蚀问题，亟待开发新型水冷一体式结构；在火焰标定方面，现有标准火焰（如NPL标准火焰）的温度上限仅为2260K，无法覆盖航空发动机和火箭发动机的标定需求。

本文定位为将工程热物理与精密光谱测量深度融合，构建一套最高工作温度可达3000K的新型标定平台。具体研究内容包括：设计带水冷一体式非预混平面绝热火焰炉结构、搭建高频TDLAS光谱测量系统、开发完整的数据处理算法、开展火焰温度标定实验验证。